\newpage
\section{Sottosistema Hardware}

\subsection{Implementazione Fisica e Tecniche di Assemblaggio SMT}

L'implementazione del prototipo ha seguito standard industriali per garantire l'affidabilità meccanica e l'integrità elettrica delle connessioni ad alta frequenza. Il montaggio dei componenti è stato eseguito tramite tecnologia \textbf{SMT}\footnote{Surface Mount Technology: tecnologia di assemblaggio che prevede il montaggio dei componenti direttamente sulla superficie del PCB.} (\textit{Surface Mount Technology}), scelta dettata da precise necessità progettuali:

\begin{itemize}
    \item \textbf{Riduzione delle Induttanze Parassite}: A differenza della tecnologia \textit{Through-Hole}, l'assenza di reofori (pin passanti) lunghi minimizza le induttanze parassite sulle linee di clock dell'RP2040 (\qty{133}{MHz}) e sui bus di comunicazione rapida come il QSPI della Flash. Questo riduce drasticamente fenomeni di \textit{ringing} e riflessioni del segnale, preservando l'integrità dei fronti d'onda digitali.
    \item \textbf{Ottimizzazione del Piano di Massa}: La tecnologia SMT permette una connessione diretta dei \textit{thermal pad} al piano di massa interno attraverso la tecnica dei \textit{via-in-pad}. Questo è fondamentale non solo per la dissipazione termica del SoC, ma anche per offrire un riferimento di tensione estremamente stabile e a bassa impedenza.
\end{itemize}



\subsubsection*{Saldatura a rifusione del modulo SAM-M8Q}

Particolare attenzione è stata dedicata al modulo GNSS \textbf{u-blox SAM-M8Q} \cite{sam_m8q_ds}, caratterizzato da un package di tipo \textbf{LCC} (\textit{Leadless Chip Carrier}). A causa della posizione dei pad disposti lungo il perimetro inferiore del componente, è stata adottata una tecnica di \textbf{saldatura a rifusione controllata} (\textit{Reflow Soldering}).



L'utilizzo di pasta saldante con flussante \textit{no-clean} e l'applicazione di un profilo termico accurato (rispettando le rampe di \textit{pre-heat} e \textit{soak} suggerite dal produttore) ha permesso di ottenere diversi vantaggi critici:
\begin{enumerate}
    \item \textbf{Prevenzione dei cortocircuiti}: La precisione del deposito della pasta saldante ha evitato la formazione di \textit{solder bridges} tra i pad adiacenti, caratterizzati da un \textit{pitch} estremamente ridotto.
    \item \textbf{Self-alignment}: È stato sfruttato il fenomeno della tensione superficiale dello stagno fuso per garantire l'auto-allineamento del modulo sulle piazzole del PCB, minimizzando gli errori di piazzamento manuale.
    \item \textbf{Integrità RF}: Una bagnabilità ottimale dei pad relativi all'antenna integrata è essenziale per minimizzare le perdite di inserzione RF. Una saldatura non perfetta in questo punto degraderebbe sensibilmente il rapporto segnale-rumore (SNR), ostacolando il fix dei satelliti in ambienti chiusi o con scarsa visibilità del cielo.
\end{enumerate}


\subsection{Alimentazione e Power Integrity}

Il design del sottosistema di alimentazione è stato ingegnerizzato per coniugare un'elevata efficienza energetica con una rigorosa pulizia dei binari di tensione, fattore critico per il corretto funzionamento di componenti a segnale misto (\textit{mixed-signal}). Il sistema accetta un ingresso nominale a 12\,V, processato attraverso una topologia di regolazione a cascata progettata per isolare il rumore di commutazione dalla logica sensibile.

\subsubsection*{Conversione Switching ad Alta Efficienza}

Il primo stadio di regolazione è deputato alla gestione dei carichi di potenza, inclusi la matrice di LED RGB e le periferiche di comunicazione, che richiedono correnti di picco elevate senza compromettere il bilancio termico del PCB.

\paragraph{Regolatore Buck: TPS563201}
Per convertire la tensione d'ingresso ai 5\,V di sistema, è stato adottato il regolatore buck sincrono \textbf{TPS563201} \cite{tps563201_ds}. Questo convertitore DC-DC è in grado di erogare correnti costanti fino a 3\,A con un'efficienza superiore al 90\%. L'architettura sincrona, grazie all'integrazione di MOSFET a bassa $R_{DS(on)}$, minimizza le perdite di conduzione. Tale scelta riduce drasticamente il calore dissipato, permettendo un design compatto senza la necessità di dissipatori passivi ingombranti. Lo schema circuitale del modulo buck è riportato in Fig. \ref{fig:buck_schematic}.

\subsubsection*{Regolazione Lineare e Protezione RF}

La sensibilità del sistema GNSS è direttamente proporzionale alla qualità dell'alimentazione fornita al modulo SAM-M8Q e ai riferimenti analogici dell'RP2040.

\paragraph{Dual-Stage Regulation: Buck + LDO}
La logica a 3.3\,V è alimentata tramite un regolatore \textbf{LDO} (\textit{Low Drop-Out}) posto in serie all'uscita dello stadio buck (Fig. \ref{fig:ldo_schematic}). Questa configurazione a due stadi funge da filtro passa-basso attivo: mentre il buck gestisce il salto di tensione principale in modo efficiente, l'LDO sfrutta il suo elevato \textbf{PSRR} (\textit{Power Supply Rejection Ratio}) per sopprimere il \textit{ripple} residuo e lo \textit{switching noise} ad alta frequenza generato dal primo stadio.

\paragraph{Integrità del segnale GNSS}
Questa strategia è vitale per prevenire interferenze elettromagnetiche (EMI) condotte che potrebbero accoppiarsi con lo stadio di ingresso RF del modulo \textbf{SAM-M8Q}. Un'alimentazione rumorosa degraderebbe il rapporto segnale-rumore (\textit{SNR}), aumentando i tempi di \textit{Time To First Fix} (TTFF) e riducendo la precisione del posizionamento \textit{indoor} o in condizioni di scarso \textit{link budget} satellitare.

\subsubsection*{Analisi della Power Integrity (PI)}

La distribuzione della potenza sul PCB è stata ottimizzata per minimizzare l'impedenza del \textit{Power Delivery Network} (PDN) e prevenire transienti di tensione localizzati.

\paragraph{Rete di Disaccoppiamento e MLCC}
Per garantire la stabilità dinamica, sono stati posizionati condensatori di bypass ceramici (\textit{MLCC}) da 100\,nF con bassa \textbf{ESR} (\textit{Equivalent Series Resistance}) in stretta prossimità di ogni pin di alimentazione del SoC. Questi componenti agiscono come serbatoi di energia a bassissima induttanza parassita, filtrando le correnti di commutazione ad alta frequenza generate dai core ARM operanti a 133\,MHz.



\paragraph{Ground Plane e Stabilità del Bus}
L'impiego di piani di massa generosi e una corretta strategia di disaccoppiamento prevengono fenomeni di \textit{ground bounce} e \textit{voltage sag}. Questa accortezza garantisce che il riferimento analogico per il modulo GPS e la stabilità del bus QSPI della Flash non vengano compromessi dai rapidi transienti di commutazione digitale, preservando l'integrità dei dati durante le operazioni di log ad alta velocità.


\subsection{Il SoC Raspberry Pi RP2040}
Il sistema orbita attorno al SoC \textbf{RP2040} \cite{rp2040_ds}, caratterizzato da un'architettura dual-core ARM Cortex-M0+ con bus AMBA AHB \cite{rp2040_ds}. La struttura simmetrica permette di parallelizzare l'acquisizione dai sensori e la scrittura su memoria flash.

\subsubsection*{Meccanismo di Interrupt vs Polling}

La scelta del paradigma di gestione degli input è critica per garantire il determinismo del sistema e la reattività dell'interfaccia utente. Si è scelto di escludere la tecnica del \textit{polling} — ovvero il campionamento ciclico dello stato dei pin all'interno del ciclo principale (\textit{super-loop}) — poiché risulterebbe intrinsecamente inefficiente. In un sistema che esegue task computazionalmente intensivi o a latenza variabile, come il parsing delle sentenze NMEA o le operazioni di scrittura su memoria Flash, il polling comporterebbe un elevato rischio di \textit{aliasing temporale}, portando alla perdita di eventi di pressione se questi occorrono mentre la CPU è occupata in altre routine.

L'adozione degli \textbf{interrupt hardware (IRQ)} sposta la gestione dell'evento dal dominio software a quello hardware, permettendo una gestione asincrona. Al verificarsi di una transizione di stato sul pin GPIO (fronte di discesa), il controller degli interrupt del core ARM Cortex-M0+ forza una sospensione immediata del programma principale per eseguire una \textit{Interrupt Service Routine} (ISR). Questo garantisce una latenza di risposta minima e costante, svincolando la reattività del sistema dal carico computazionale istantaneo.



\paragraph{Architettura degli Interrupt sull'RP2040}
L'efficienza di questa implementazione è legata alla specifica architettura del SoC RP2040. Come documentato nel datasheet \cite[\href{https://datasheets.raspberrypi.com/rp2040/rp2040-datasheet.pdf\#page=240}{Sez. 2.19.3}]{rp2040_ds}, il sistema non assegna un vettore di interrupt dedicato ad ogni singolo pin, ma aggrega le richieste tramite una logica \textbf{OR-wired} per ogni banco di memoria GPIO. 

Dato che tutti i pulsanti del dispositivo afferiscono allo stesso banco, la logica di gestione è stata strutturata in due fasi logiche:
\begin{enumerate}
    \item \textbf{Segnalazione Hardware:} L'evento sul pin solleva un segnale di interrupt globale verso il processore.
    \item \textbf{Discriminazione Software:} All'interno della ISR, il software interroga i registri di stato degli interrupt (\textit{Interrupt Status Registers}) per identificare quale specifico GPIO abbia scatenato l'evento. Questa architettura centralizzata permette di ottimizzare l'uso delle risorse, utilizzando un unico handler per gestire l'intero set di input utente.
\end{enumerate}

\paragraph{Condizionamento e Pull-up}
Per garantire la stabilità dei livelli logici, i pin sono stati configurati in modalità \texttt{INPUT\_PULLUP} tramite il framework Arduino. Questa impostazione attiva una rete di resistenze interne al SoC (con valori nominali compresi tra $50\,k\Omega$ e $80\,k\Omega$). Tale configurazione, lavorando in sinergia con il filtro RC esterno visibile nello schematico di Fig. \ref{fig:button_sch}, assicura un pull-up ``forte'' che previene stati di alta impedenza (\textit{floating}) e mitiga l'effetto di accoppiamenti capacitivi indesiderati, garantendo una lettura pulita anche in ambienti elettricamente rumorosi.